\documentclass[10pt,a4paper]{amsart}
\usepackage{amsmath,amssymb,amsthm,graphicx}
% Retained minimal preamble for the Stacks Project chapter flat.tex (More on Flatness).
% Verbatim excerpt of the chapter preamble: Xy-pic setup (native preamble lines 14-18) and the native
% theorem declarations with the mathematical macros (native preamble lines 154-217). Copied mechanically;
% nothing added, removed or redefined. The wrapper class and the amsmath/amssymb/amsthm/graphicx base
% packages are supplied by the assembler, so the class/externaldocument part of the native preamble is omitted.
% For commutative diagrams we use Xy-pic
\usepackage[all]{xy}

% We use 2cell for 2-commutative diagrams.
\xyoption{2cell}
% Theorem environments.
%
\theoremstyle{plain}
\newtheorem{theorem}[subsection]{Theorem}
\newtheorem{proposition}[subsection]{Proposition}
\newtheorem{lemma}[subsection]{Lemma}

\theoremstyle{definition}
\newtheorem{definition}[subsection]{Definition}
\newtheorem{example}[subsection]{Example}
\newtheorem{exercise}[subsection]{Exercise}
\newtheorem{situation}[subsection]{Situation}

\theoremstyle{remark}
\newtheorem{remark}[subsection]{Remark}
\newtheorem{remarks}[subsection]{Remarks}

\numberwithin{equation}{subsection}

% Macros
%
\def\lim{\mathop{\mathrm{lim}}\nolimits}
\def\colim{\mathop{\mathrm{colim}}\nolimits}
\def\Spec{\mathop{\mathrm{Spec}}}
\def\Hom{\mathop{\mathrm{Hom}}\nolimits}
\def\Ext{\mathop{\mathrm{Ext}}\nolimits}
\def\SheafHom{\mathop{\mathcal{H}\!\mathit{om}}\nolimits}
\def\SheafExt{\mathop{\mathcal{E}\!\mathit{xt}}\nolimits}
\def\Sch{\mathit{Sch}}
\def\Mor{\mathop{\mathrm{Mor}}\nolimits}
\def\Ob{\mathop{\mathrm{Ob}}\nolimits}
\def\Sh{\mathop{\mathit{Sh}}\nolimits}
\def\NL{\mathop{N\!L}\nolimits}
\def\CH{\mathop{\mathrm{CH}}\nolimits}
\def\proetale{{pro\text{-}\acute{e}tale}}
\def\etale{{\acute{e}tale}}
\def\QCoh{\mathit{QCoh}}
\def\Ker{\mathop{\mathrm{Ker}}}
\def\Im{\mathop{\mathrm{Im}}}
\def\Coker{\mathop{\mathrm{Coker}}}
\def\Coim{\mathop{\mathrm{Coim}}}

% Boxtimes
%
\DeclareMathSymbol{\boxtimes}{\mathbin}{AMSa}{"02}

%
% Macros for moduli stacks/spaces
%
\def\QCohstack{\mathcal{QC}\!\mathit{oh}}
\def\Cohstack{\mathcal{C}\!\mathit{oh}}
\def\Spacesstack{\mathcal{S}\!\mathit{paces}}
\def\Quotfunctor{\mathrm{Quot}}
\def\Hilbfunctor{\mathrm{Hilb}}
\def\Curvesstack{\mathcal{C}\!\mathit{urves}}
\def\Polarizedstack{\mathcal{P}\!\mathit{olarized}}
\def\Complexesstack{\mathcal{C}\!\mathit{omplexes}}
% \Pic is the operator that assigns to X its picard group, usage \Pic(X)
% \Picardstack_{X/B} denotes the Picard stack of X over B
% \Picardfunctor_{X/B} denotes the Picard functor of X over B
\def\Pic{\mathop{\mathrm{Pic}}\nolimits}
\def\Picardstack{\mathcal{P}\!\mathit{ic}}
\def\Picardfunctor{\mathrm{Pic}}
\def\Deformationcategory{\mathcal{D}\!\mathit{ef}}

% --- appended 2026-10-12 (Chief): the project own definitions of the reference/slogan/history
% display environments, copied mechanically from the shared Stacks preamble lines 8-12
% (\usepackage{verbatim} plus the three \newenvironment lines, which the chapter preamble keeps
% verbatim from the project preamble). Nothing added, removed or redefined: these environments
% carry no mathematical content and the project itself typesets them as comments. Without these
% three lines a counted statement whose own text uses them fails with Environment slogan undefined.
\usepackage{verbatim}
\newenvironment{reference}{\comment}{\endcomment}
\newenvironment{slogan}{\comment}{\endcomment}
\newenvironment{history}{\comment}{\endcomment}
\usepackage{cleveref}
\crefname{originallemma}{original lemma}{original lemmas}
\Crefname{originallemma}{Original lemma}{Original lemmas}
\makeatletter
\newlabel{morphisms-lemma-closed-immersion-finite-presentation}{{01TV}{source}{}{sourcepointer-1}{}}
\newlabel{morphisms-lemma-closed-immersion-finite-presentation@cref}{{[originallemma][0][]01TV}{[1][0][]source}{}{}{}}
\makeatother
\crefname{originallemma}{original lemma}{original lemmas}
\Crefname{originallemma}{Original lemma}{Original lemmas}
\makeatletter
\newlabel{divisors-lemma-separate-disjoint-opens-by-blowing-up}{{080P}{source}{}{sourcepointer-2}{}}
\newlabel{divisors-lemma-separate-disjoint-opens-by-blowing-up@cref}{{[originallemma][0][]080P}{[1][0][]source}{}{}{}}
\makeatother
\crefname{originaldefinition}{original definition}{original definitions}
\Crefname{originaldefinition}{Original definition}{Original definitions}
\makeatletter
\newlabel{divisors-definition-admissible-blowup}{{080K}{source}{}{sourcepointer-3}{}}
\newlabel{divisors-definition-admissible-blowup@cref}{{[originaldefinition][0][]080K}{[1][0][]source}{}{}{}}
\makeatother
\crefname{originallemma}{original lemma}{original lemmas}
\Crefname{originallemma}{Original lemma}{Original lemmas}
\makeatletter
\newlabel{morphisms-lemma-composition-finite-presentation}{{01TR}{source}{}{sourcepointer-4}{}}
\newlabel{morphisms-lemma-composition-finite-presentation@cref}{{[originallemma][0][]01TR}{[1][0][]source}{}{}{}}
\makeatother
\crefname{originallemma}{original lemma}{original lemmas}
\Crefname{originallemma}{Original lemma}{Original lemmas}
\makeatletter
\newlabel{divisors-lemma-strict-transform}{{080E}{source}{}{sourcepointer-5}{}}
\newlabel{divisors-lemma-strict-transform@cref}{{[originallemma][0][]080E}{[1][0][]source}{}{}{}}
\makeatother
\crefname{originallemma}{original lemma}{original lemmas}
\Crefname{originallemma}{Original lemma}{Original lemmas}
\makeatletter
\newlabel{properties-lemma-quasi-coherent-finite-type-ideals}{{01PH}{source}{}{sourcepointer-6}{}}
\newlabel{properties-lemma-quasi-coherent-finite-type-ideals@cref}{{[originallemma][0][]01PH}{[1][0][]source}{}{}{}}
\makeatother
\crefname{originallemma}{original lemma}{original lemmas}
\Crefname{originallemma}{Original lemma}{Original lemmas}
\makeatletter
\newlabel{properties-lemma-extend}{{01PF}{source}{}{sourcepointer-7}{}}
\newlabel{properties-lemma-extend@cref}{{[originallemma][0][]01PF}{[1][0][]source}{}{}{}}
\makeatother
\title{Two locally closed pieces that are disjoint over a quasi-compact open can be separated by an admissible blowup}
\author{The Stacks Project Authors}
\date{stacks-project:flat.tex:local-snapshot-20261011; GNU FDL 1.2 or later}
\begin{document}
\maketitle
\noindent\textit{Editorial note.} All delivered mathematical text is the author's own, copied from the pinned Stacks Project chapter file flat.tex; no statement or proof step was rewritten, reordered or removed.\par\medskip
\noindent\textit{Editorial note.} The preamble is the chapter's own retained minimal preamble (sources/stacks-flat/retained-preamble.tex), a verbatim excerpt of the native Xy-pic setup and of the native theorem declarations and mathematical macros, followed by the project own definitions of the reference, slogan and history display environments (verbatim package plus three newenvironment lines) that the retained file carries at its end and that the project itself typesets as comments. No macro was invented or redefined, and the wrapper class and the amsmath/amssymb/amsthm/graphicx base packages are supplied by the assembler.\par\medskip
\noindent\textit{Editorial note.} The counted statement and proof contain no cross reference, no citation and no included file, so the delivered body is self contained.\par\medskip
\noindent\textit{Editorial note.} Canonical location of the counted proposition: Stacks tag 0F3V (sources/stacks-flat/pinned-original-source.tex line 9554).\par\medskip
\noindent\textit{Allowed standard prerequisites (editorial).}\begin{itemize}
\item Finite presentation of a closed subscheme and finite type of its ideal sheaf. Morphisms, Lemma labelled morphisms-lemma-closed-immersion-finite-presentation, cited at proof lines 9568-9572. Allows the ideals I\_i to be used in a blowup construction.
\item Separating disjoint closed subschemes over an open by an admissible blowup. Divisors, Lemma labelled divisors-lemma-separate-disjoint-opens-by-blowing-up, cited at proof lines 9574-9579. Produces the blowup of Z making the strict transforms disjoint.
\item Strict transforms under a blowup and the composition of finite presentation morphisms. Divisors, Lemma labelled divisors-lemma-strict-transform, cited at proof lines 9594-9598, and Morphisms, Lemma labelled morphisms-lemma-composition-finite-presentation, cited at proof lines 9587-9591. Shows the blowup of the center Y on X is admissible and computes the strict transforms.
\item Constructible closed subsets as zero loci of finite type ideals, and extension of ideals from an open. Properties, Lemmas labelled properties-lemma-quasi-coherent-finite-type-ideals and properties-lemma-extend, cited at proof lines 9595-9603. Reduces part (2) to part (1).
\end{itemize}\medskip
\ifdefined\hypertarget\hypertarget{sourcepointer-1}{}\fi
\noindent\textit{Source reference.} Original lemma 01TV: sources/stacks-chow/tags.txt, source lines 2388--2388. Cross-chapter standard prerequisite located in the project tag list staged in the local snapshot; invoked at proof line 9572 of the counted proof. It is retained only as a navigation pointer; every load-bearing step is inside the delivered span.\par\medskip
\ifdefined\hypertarget\hypertarget{sourcepointer-2}{}\fi
\noindent\textit{Source reference.} Original lemma 080P: sources/stacks-chow/tags.txt, source lines 10189--10189. Cross-chapter standard prerequisite located in the project tag list staged in the local snapshot; invoked at proof line 9577 of the counted proof. It is retained only as a navigation pointer; every load-bearing step is inside the delivered span.\par\medskip
\ifdefined\hypertarget\hypertarget{sourcepointer-3}{}\fi
\noindent\textit{Source reference.} Original definition 080K: sources/stacks-chow/tags.txt, source lines 10185--10185. Cross-chapter standard prerequisite located in the project tag list staged in the local snapshot; invoked at proof line 9583 of the counted proof. It is retained only as a navigation pointer; every load-bearing step is inside the delivered span.\par\medskip
\ifdefined\hypertarget\hypertarget{sourcepointer-4}{}\fi
\noindent\textit{Source reference.} Original lemma 01TR: sources/stacks-chow/tags.txt, source lines 2384--2384. Cross-chapter standard prerequisite located in the project tag list staged in the local snapshot; invoked at proof line 9584 of the counted proof. It is retained only as a navigation pointer; every load-bearing step is inside the delivered span.\par\medskip
\ifdefined\hypertarget\hypertarget{sourcepointer-5}{}\fi
\noindent\textit{Source reference.} Original lemma 080E: sources/stacks-chow/tags.txt, source lines 10179--10179. Cross-chapter standard prerequisite located in the project tag list staged in the local snapshot; invoked at proof line 9590 of the counted proof. It is retained only as a navigation pointer; every load-bearing step is inside the delivered span.\par\medskip
\ifdefined\hypertarget\hypertarget{sourcepointer-6}{}\fi
\noindent\textit{Source reference.} Original lemma 01PH: sources/stacks-chow/tags.txt, source lines 2230--2230. Cross-chapter standard prerequisite located in the project tag list staged in the local snapshot; invoked at proof line 9595 of the counted proof. It is retained only as a navigation pointer; every load-bearing step is inside the delivered span.\par\medskip
\ifdefined\hypertarget\hypertarget{sourcepointer-7}{}\fi
\noindent\textit{Source reference.} Original lemma 01PF: sources/stacks-chow/tags.txt, source lines 2228--2228. Cross-chapter standard prerequisite located in the project tag list staged in the local snapshot; invoked at proof line 9599 of the counted proof. It is retained only as a navigation pointer; every load-bearing step is inside the delivered span.\par\medskip
\begin{lemma}
\label{lemma-separate-disjoint-locally-closed-by-blowing-up}
Let $X$ be a quasi-compact and quasi-separated scheme.
Let $U \subset X$ be a quasi-compact open.
\begin{enumerate}
\item If $Z_1, Z_2 \subset X$ are closed subschemes of finite
presentation such that $Z_1 \cap Z_2 \cap U = \emptyset$, then
there exists a $U$-admissible blowing up $X' \to X$
such that the strict transforms of $Z_1$ and $Z_2$ are disjoint.
\item If $T_1, T_2 \subset U$ are disjoint constructible closed subsets, then
there is a $U$-admissible blowing up $X' \to X$ such that the closures of
$T_1$ and $T_2$ are disjoint.
\end{enumerate}
\end{lemma}

\begin{proof}
Proof of (1). The assumption that $Z_i \to X$ is of finite presentation
signifies that the quasi-coherent ideal sheaf $\mathcal{I}_i$ of $Z_i$
is of finite type, see 
Morphisms, Lemma \ref{morphisms-lemma-closed-immersion-finite-presentation}.
Denote $Z \subset X$ the closed subscheme
cut out by the product $\mathcal{I}_1 \mathcal{I}_2$.
Observe that $Z \cap U$ is the disjoint union
of $Z_1 \cap U$ and $Z_2 \cap U$. By Divisors, Lemma
\ref{divisors-lemma-separate-disjoint-opens-by-blowing-up}
there is a $U \cap Z$-admissible blowup $Z' \to Z$ such that
the strict transforms of $Z_1$ and $Z_2$ are disjoint.
Denote $Y \subset Z$ the center of this blowing up.
Then $Y \to X$ is a closed immersion of finite presentation as the composition
of $Y \to Z$ and $Z \to X$ (Divisors, Definition
\ref{divisors-definition-admissible-blowup}
and Morphisms, Lemma \ref{morphisms-lemma-composition-finite-presentation}).
Thus the blowing up $X' \to X$ of $Y$ is a $U$-admissible blowing
up. By general properties of strict transforms, the
strict transform of $Z_1, Z_2$ with respect to $X' \to X$
is the same as the strict transform of $Z_1, Z_2$ with respect
to $Z' \to Z$, see
Divisors, Lemma \ref{divisors-lemma-strict-transform}.
Thus (1) is proved.

\medskip\noindent
Proof of (2). By Properties, Lemma
\ref{properties-lemma-quasi-coherent-finite-type-ideals}
there exists a finite type quasi-coherent sheaf of ideals
$\mathcal{J}_i \subset \mathcal{O}_U$ such that
$T_i = V(\mathcal{J}_i)$ (set theoretically).
By Properties, Lemma \ref{properties-lemma-extend}
there exists a finite type quasi-coherent sheaf
of ideals $\mathcal{I}_i \subset \mathcal{O}_X$
whose restriction to $U$ is $\mathcal{J}_i$.
Apply the result of part (1) to the closed
subschemes $Z_i = V(\mathcal{I}_i)$ to conclude.
\end{proof}

\end{document}
